\begin{textsample}{NEG Dim 3 – 2010\_c6 – Score -7.00 – 2010\_c6/enade\_2010\_tec\_agro\_ind\_8fdf27ddf8.txt}  \label{ex:f3_neg_012}
CAPÍTULO 3 ANÁLISE TÉCNICA DA PROVA Este capítulo tem por objetivo apresentar o desempenho dos estudantes de Tecnologia em Agroindústria no ENADE/2010.
Para isso, foram calculadas as estatísticas básicas da prova em seu todo, bem como as estatísticas dos componentes relacionadas à Formação Geral e ao de Conhecimento Específico da Área.
Nas tabelas, são apresentados o tamanho da população inscrita e de presentes, e as seguintes estatísticas das notas7 : média do desempenho na prova, erro padrão da média, desvio padrão, nota mínima, mediana, nota máxima, coeficiente de assimetria e de variação.
As estatísticas apresentadas neste capítulo contemplam o total de estudantes e, separadamente, os Ingressantes e os Concluintes.
Foram calculadas tendo -se em vista as seguintes agregações : ( a ) as Grandes Regiões e o país como um todo ; ( b ) a Categoria Administrativa ; e ( c ) a Organização Acadêmica.
Em relação aos gráficos de distribuição de notas, o intervalo considerado foi de 10 unidades, aberto à esquerda e fechado à direita, com exceção do primeiro intervalo, [ 0 ; 10 ], fechado em ambos os extremos.
Para os gráficos de distribuição das notas das questões discursivas, foram consideradas mais duas categorias : questão em branco e nota zero. 3.1 ESTATÍSTICAS BÁSICAS DA PROVA 3.1.1 Estatísticas Básicas Gerais A Tabela 3.1 apresenta as estatísticas básicas da prova por grupo de estudantes.
A população total de estudantes regularmente inscritos8 para a prova foi de 387.
Destes, 287 estiveram presentes, sendo de 25,8 \% o índice de não comparecimento.
A abstenção foi maior entre Ingressantes ( 29,0 \% ) do que entre Concluintes ( 16,5 \% ).
A nota média de toda a prova foi 43,6, sendo que os Ingressantes obtiveram média mais baixa, de 42,1, que os Concluintes, cuja média foi 47,4.
O desvio padrão foi 13,4, sendo o do grupo de Ingressantes menor ( 12,7 ) do que o do grupo dos Concluintes ( 14,4 ) ; os coeficientes de variação foram 30,2 \% no grupo dos Ingressantes e 30,4 \% no dos Concluintes.
A nota máxima foi 74,5, obtida por Concluintes, ao passo que a maior nota obtida por Ingressantes foi 73,9. 7 Algumas dessas estatísticas estão definidas no Capítulo 1. 8 As estatísticas das notas apresentadas consideram apenas os estudantes regularmente inscritos e com resultados validados por ausência de erro administrativo.
Uma comparação entre o desempenho de Ingressantes e Concluintes também pode ser realizada a partir dos dados apresentados no Gráfico 3.1, com a distribuição de notas finais.
Destaca -se que as notas mais elevadas, considerando -se o intervalo posterior à nota 40,0, foram atingidas principalmente pelos Concluintes, cuja concentração das notas está à direita dos Ingressantes.
Já nas faixas de notas mais baixas, inverteu -se o quadro identificado nas faixas de notas mais altas.
Essa característica é corroborada pelo valor da nota mediana 43,1 no grupo dos Ingressantes, abaixo da encontrada no grupo dos Concluintes, 49,0, transparecendo superioridade de desempenho dos Concluintes em relação aos Ingressantes na prova em tela ( Tabela 3.1 ).
A faixa modal para Ingressantes foi a que vai de 40 a 50 pontos com 35,9 \% dos estudantes do grupo e para Concluintes, a faixa seguinte, de 50 a 60 pontos com 27,2 \%.
Os coeficientes de assimetria indicam uma assimetria negativa para os Concluintes ( -0,6 ) e para Ingressantes ( -0,1 ).
O Gráfico 3.2 apresenta informações referentes à nota final de Ingressantes e Concluintes, comparando -se os resultados de acordo com as Grandes Regiões do país, a Categoria Administrativa e a Organização Acadêmica.
Considerando -se as notas médias dos alunos em cada Grande Região, ressalta -se que não houve estudantes, Ingressantes ou Concluintes, participando do exame na região Sudeste.
Especificamente quanto aos Concluintes, só houve participantes nas regiões Norte e Sul.
Em relação aos Ingressantes, a única região em que não houve representantes foi a Sudeste, como já mencionado.
Em relação às notas aferidas por região em âmbito nacional ( média 42,1 para Ingressantes e 47,4 para Concluintes ), não existe uma diferença estatisticamente significativa ao nível de 95 \%, para Ingressantes, entre a menor média, obtida na região Centro-Oeste ( 39,3 ) e a maior, obtida na região Norte ( 43,2 ).
Entre as duas regiões onde houve Concluintes, a maior nota média foi registrada na região Norte, com 52,4 ( média superior à nacional ), enquanto a região Sul obteve 41,3.
Por outro lado, esta diferença entre Concluintes foi estatisticamente significativa ao nível de 95 \% Levando -se em conta os grupos dos estudantes em Categorias Administrativas, só houve a participação de instituições públicas.
Nelas, Concluintes tiveram nota média 47,4, enquanto Ingressantes atingiram 42,1.
Destaca -se que ambas as médias são idênticas à nacional.
No que diz respeito às Organizações Acadêmicas, não houve participação de estudantes em Centros Universitários e Faculdades.
Nas Universidades, Concluintes obtiveram nota média 47,4 e Ingressantes, 42,1.
Conclui -se que todos os estudantes que realizaram o exame, na Área de Tecnologia em Agroindústria, portanto, eram de Universidades públicas. 3.1.2 Estatísticas Básicas no Componente de Formação Geral A Tabela 3.2 apresenta as estatísticas básicas em relação ao componente da prova que avalia a Formação Geral dos estudantes.
Os Concluintes obtiveram desempenho médio 47,5, ligeiramente inferior ao dos Ingressantes, 47,8.
Quanto à variabilidade das notas nos dois grupos, o desvio padrão das notas obtidas esteve em torno de 16,5.
Os dois grupos alcançaram notas máximas aproximadas, sendo 88,0 a maior nota obtida por pelo menos um Ingressante, e 84,0 a maior nota entre os Concluintes.
O Gráfico 3.3 propicia a comparação entre Concluintes e Ingressantes relativa ao desempenho no componente de Formação Geral.
Dos Ingressantes, 24,8 \% obtiveram notas na faixa entre 50 e 60 pontos, sendo essa a faixa modal.
Nesse mesmo intervalo, Concluintes obtiveram 24,7 \%.
Nas faixas de notas seguintes, Concluintes e Ingressantes se alternam em rendimentos mais elevados.
No intervalo entre 60 e 70 pontos, por exemplo, entre Ingressantes, houve 12,1 \% dos participantes, ao passo que, entre Concluintes, 13,6 \%.
No intervalo seguinte, de 70 a 80 pontos, verificou -se a incidência de 7,8 \% de Ingressantes e 7,4 \% de Concluintes.
A faixa modal para Concluintes foi, entretanto, entre 40 e 50 pontos com 25,9 \% deste grupo.
Apenas 2,5 \% dos Concluintes tiveram desempenho na menor faixa de notas, que vai de zero a 10 pontos.
O coeficiente de assimetria foi negativo para Concluintes ( -0,4 ) e para Ingressantes ( -0,1 ).
No Gráfico 3.4 são apresentadas as informações referentes ao desempenho dos Ingressantes e dos Concluintes no componente de Formação Geral, em diferentes agregações como Grande Região do país, Categoria Administrativa e Organização Acadêmica.
Quanto às notas médias segundo cada Grande Região, deve -se considerar as peculiaridades do curso de Tecnologia em Agroindústria, cuja participação de estudantes se deu, no caso de Concluintes, apenas nas regiões Norte e Sul.
Em relação a Ingressantes, a única região sem participantes foi a Sudeste.
No grupo de Ingressantes não houve diferença estatisticamente significativa ao nível de 95 \% entre as médias nas regiões Nordeste ( 53,8, a mais alta ) e Sul ( 46,5, a mais baixa ).
No grupo de Concluintes, as notas foram 53,6 na região Norte e 40,2, na região Sul e a diferença também não foi estatisticamente significativa ao nível de 95 \%.
Tendo como foco as Categorias Administrativas, houve participação, conforme já fora informado, somente em instituições públicas, nas quais Ingressantes obtiveram nota média 47,8, e Concluintes, 47,5.
Quanto às Organizações Acadêmicas, só houve participação de Universidades.
Ingressantes alcançaram nota 47,8 e Concluintes, 47,5.
Observa -se que nas Universidades Públicas, portanto, Ingressantes tiveram desempenho ligeiramente superior ao desempenho dos Concluintes. 3.1.3 Estatísticas Básicas do Componente de Conhecimento Específico A Tabela 3.3 apresenta as estatísticas básicas referentes ao componente de Conhecimento Específico da Área de Tecnologia em Agroindústria.
A média do desempenho foi 42,2.
Os Concluintes apresentaram desempenho médio ( 47,3 ) superior ao dos Ingressantes ( 40,2 ) ; as medidas de dispersão – desvio padrão e amplitude geral das notas – são menores no grupo dos Ingressantes do que no grupo dos Concluintes.
Os coeficientes de variação, 35,1 \% no grupo dos Ingressantes e 34,7 \% no grupo dos Concluintes, indicam dispersão relativa menor nas notas desses últimos.
As notas máximas obtidas por Concluintes e Ingressantes foram 77,8 e 80,8, respectivamente, notas inferiores àquelas encontradas no componente de Formação Geral ( 84,0 e 88,0, respectivamente ).
As medianas das distribuições das notas dos Ingressantes e dos Concluintes foram 40,3 e 49,2, respectivamente, enfatizando a superioridade dos Concluintes neste componente da prova.
Assim como os Gráficos 3.1 e 3.3, o Gráfico 3.5, apresentado a seguir, proporciona uma comparação do desempenho entre Concluintes e Ingressantes, em relação ao componente de Conhecimento Específico.
Dentre as três comparações feitas, observa -se que, embora a distribuição das notas seja semelhante à apresentada na prova como um todo, neste componente, a diferença de desempenho entre Concluintes e Ingressantes é mais acentuada, destacando -se o desempenho mais elevado dos primeiros.
Claramente a distribuição das notas dos Ingressantes, mais concentradas à esquerda em relação aos Concluintes, evidencia o percentual mais alto de notas nas faixas mais baixas para este grupo.
A maior concentração de Concluintes está no intervalo de ( 50 a 60 ] pontos com 24,7 \% dos estudantes deste grupo, ao passo que a de Ingressantes está em ( 40 a 50 ] pontos com 28,2 \% destes.
Assim como os Gráficos 3.2 e 3.4, o Gráfico 3.6 apresenta uma comparação dos resultados em relação às Grandes Regiões do país, às Categorias Administrativas e às Organizações Acadêmicas, agora levando em conta o desempenho dos Ingressantes e dos Concluintes no componente de Conhecimento Específico da prova.
Ao se observar as notas médias dos estudantes, segundo cada Grande Região do país, nas duas regiões em que houve participação de Concluintes, a média mais elevada foi encontrada na região Norte ( 52,0 ), maior que a média nacional ( 47,3 ).
A região Sul obteve média 41,7.
Em relação aos Ingressantes, também a região Norte apresentou a média mais elevada ( 41,6 ), seguida da região Sul, com média igual a 40,4.
A média mais baixa entre os Ingressantes foi verificada na região Centro-Oeste ( 36,7 ).
Quanto à Categoria Administrativa, havendo participação apenas em instituições públicas, Concluintes obtiveram nota média 47,3, enquanto Ingressantes, 40,2.
Essas instituições públicas eram somente Universidades, nas quais as essas mesmas médias de notas foram obtidas, respectivamente, pelos dois grupos. 3.2 ANÁLISE DAS QUESTÕES OBJETIVAS 3.2.1 Componente de Formação Geral A Tabela 3.4 apresenta as estatísticas básicas relativas às oito questões objetivas do componente da prova que abrange a Formação Geral dos estudantes.
Como pode ser observado, os estudantes Concluintes obtiveram desempenho superior, com nota média 54,7, em relação aos alunos Ingressantes, cuja nota média foi 52,6.
As medidas de dispersão nos dois grupos de alunos foram próximas, como indicam os desvios padrões, assim como a amplitude total das notas, sendo que a nota mínima foi zero e a máxima, 100,0, para os dois grupos de alunos avaliados.
Os coeficientes de variação foram de 37,5 \% no grupo dos Ingressantes a 37,8 \% no dos Concluintes.
A mediana foi igual a 57,1 em ambos os grupos.
O coeficiente de assimetria do total de estudantes revelou assimetria negativa ( -0,1 ).
A Tabela 3.5 apresenta as classificações das questões objetivas do componente de Formação Geral, segundo o índice de facilidade.
As questões objetivas da prova foram assim avaliadas : das oito questões, nenhuma teve o índice de facilidade classificado como muito fácil.
Duas questões foram tidas como fáceis, por terem índice de acertos situado na faixa entre 0,61 e 0,85 ( de 61,0 \% a 85,0 \% de acertos ).
Três questões foram consideradas de dificuldade média, situando -se no intervalo entre 0,41 e 0,60 do índice de facilidade, ou seja, houve entre 41,0 \% e 60,0 \% de acertos.
Duas questões foram classificadas como difíceis, situando -se no intervalo entre 0,16 e 0,40.
Por fim, uma questão apresentou menos de 0,15 de acertos, razão pela qual foi classificada como muito difícil.
A Tabela 3.6 descreve os resultados das análises das questões objetivas relativas à Formação Geral, segundo o poder de discriminação, utilizando -se para tal, do índice de discriminação ponto bisserial.
Cinco das oito questões apresentaram índices acima de 0,40 e, assim, foram classificadas como muito boas nessa característica para esse grupo de alunos ; duas questões tiveram um bom índice de discriminação, entre 0,30 e 0,39.
Nenhuma das questões foi classificada como média.
A questão de número 6 teve fraco nível de discriminação.
Na sequência, na Tabela 3.7, estão especificadas, por questão, os resultados da análise relativa aos índices de facilidade e de discriminação de cada questão no componente relativo à Formação Geral da prova de Tecnologia em Agroindústria.
O índice de facilidade variou de 0,11 a 0,82, e o de discriminação, de 0,11 a 0,56.
As questões com índices de discriminação muito bom, de números 1, 3, 4, 5 e 8, figuraram entre as mais fáceis desse conjunto, sendo que as questões 3 e 8 foram classificadas como fáceis com respeito ao índice de facilidade e as demais ( 1, 4 e 5 ) como de dificuldade média.
A questão 8 foi uma das que apresentou maior poder discriminatório, com índice de 0,47, e foi também a mais fácil, com proporção de 0,82 de acertos.
A questão 6 apresentou índice de facilidade de 0,11, ou seja, um quantitativo mais reduzido de estudantes conseguiu resolvê -la, dentro do universo de participantes.
Seu coeficiente de discriminação foi também 0,11, classificado como fraco e foi então eliminada do cômputo da nota final.
O Gráfico 3.7 analisa o comportamento da questão de número 8 de Formação Geral.
Trata -se da questão mais fácil e com o maior índice de discriminação desse componente da prova.
No gráfico a seguir, cada uma das cinco curvas representa o percentual de respostas em determinada alternativa da questão, em função da nota dos estudantes nesta parte da prova ( Formação Geral/Múltipla Escolha ) antes de possíveis eliminações pelo critério do ponto bisserial.
A curva em rosa corresponde à alternativa B, a correta para este item.
Assim, observa -se que entre os estudantes com notas mais baixas, indicando desempenho mais fraco neste componente do exame, a situação mais frequente foi a escolha de uma das alternativas incorretas : a alternativa D ( em vermelho ), C ( em verde ) ou E ( em roxo ).
Na medida em que a nota aumenta, representando desempenho melhor, aumenta, concomitantemente, a proporção de estudantes que selecionaram a alternativa correta B, atingindo 100 \% para as notas acima de 80,0.
Essa análise permite verificar como a questão discriminou entre os grupos de desempenho, justificando o alto índice obtido na questão.
Os gráficos relativos às demais questões de Formação Geral constam do Anexo I. 3.2.2 Componente de Conhecimento Específico A Tabela 3.8 apresenta as estatísticas básicas em relação às questões objetivas do componente de Conhecimento Específico da prova.
Nesse conjunto de questões, os estudantes Concluintes ( 52,3 ) obtiveram nota média superior aos Ingressantes ( 45,3 ).
Quanto à variabilidade das notas, o desvio padrão no grupo dos Ingressantes foi 15,6, e no grupo dos Concluintes, 18,2.
Os coeficientes de variação foram 34,4 \% e 34,8 \%, respectivamente.
Com relação à mediana, no grupo dos Ingressantes, esta foi igual a 44,8, estando pouco abaixo da média desse grupo, enquanto no grupo dos Concluintes, a mediana foi igual a 52,6, superando um pouco a média desse grupo.
O coeficiente de assimetria foi comportamento simétrico ( 0,0 ) para Ingressantes e negativo ( -0,4 ) para Concluintes.
A Tabela 3.9 apresenta as classificações das questões objetivas da parte da prova relativa ao componente de Conhecimento Específico, segundo o índice de facilidade.
Deve -se considerar o fato de as questões de números 23 e 29 terem sido anuladas pela Comissão, de modo que a classificação foi estabelecida com base em 25 das 27 questões.
A partir dos índices obtidos, pode -se concluir que quase a metade das questões objetivas da prova foram tidas como pelo menos difíceis : das 25 questões, onze foram classificadas como difíceis e uma como muito difícil.
Não houve questões classificadas como muito fáceis, ao passo que quatro foram tidas como fáceis, na faixa de 0,61 a 0,85 do índice de facilidade, e nove consideradas médias, entre 0,41 e 0,60.
Na Tabela 3.10 são apresentados os índices de discriminação das questões objetivas do componente de Conhecimento Específico da prova.
Observa -se, segundo os índices estabelecidos, que a capacidade de discriminação das questões objetivas desse componente da prova foi classificada, para cinco das 25 questões válidas, como boa, enquanto outras cinco tiveram índice de discriminação muito bom.
Portanto, 10 em 25 questões proporcionaram índices de discriminação bons ou muito bons.
Nove delas foram classificadas como médias e seis como fracas, sendo quinze, por conseguinte, a quantidade de questões nos dois patamares mais baixos de discriminação.
Constata -se, assim, que a prova – no que se refere ao componente de Conhecimento Específico – possuía capacidade de média a fraca para discriminar, entre aqueles que dominam ou não o conteúdo.
Na Tabela 3.11 estão especificados, por questão, os resultados da análise relativa aos índices de facilidade e de discriminação de cada uma das questões do componente de Conhecimento Específico da prova de Tecnologia em Agroindústria.
Dentre as questões que alcançaram os maiores índices de discriminação, muito bom, as de números 11, 12, 14 e 37 foram classificadas como de dificuldade média, e a questão 19 como difícil.
Em termos de proporção, tais questões variaram de 0,28 a 0,58 quanto aos acertos.
Seus índices de discriminação situaram -se de 0,41 a 0,49.
A questão de número 22 foi a mais difícil dentre as 25 questões específicas válidas, com baixo índice de facilidade, de apenas 0,08, o que, em termos percentuais, corresponde a 8,0 \% de acertos.
Essa questão apresentou poder discriminatório igualmente baixo, 0,05, o que comprova ter sido esta a mais difícil para os estudantes.
As seis questões com índice de discriminação fraco ( 15, 16, 20, 22, 25 e 28 ) foram eliminadas do cômputo da nota final.
O Gráfico 3.8 analisa o comportamento da questão 14 do componente de Conhecimento Específico.
Esta foi uma das questões mais fáceis da prova, apresentando índice de facilidade de 0,58, ou seja, 58,0 \% dos estudantes assinalaram acertadamente a opção C, correspondente ao gabarito.
Seu índice de discriminação foi igual a 0,49, classificado como muito bom.
Neste gráfico, cada uma das cinco curvas representa o percentual de respostas em determinada alternativa da questão 14, em função da nota dos estudantes nesta parte da prova antes de possíveis eliminações de questões pelo critério do ponto bisserial.
A alternativa correta C, representada no gráfico pela curva em verde, foi escolhida em maiores proporções pelos alunos com desempenho melhor nesta componente do exame, enquanto as alternativas incorretas, também denominadas distratores, foram selecionadas principalmente por aqueles com notas mais baixas, ainda que aqueles com nota zero, na sua quase totalidade deixaram em branco ou marcaram mais de uma alternativa, comportamento considerado inválido.
A proporção de alunos que selecionou a resposta correta C aumenta gradativamente, chegando a atingir 100 \% para notas acima de 70,0, enquanto a proporção dos que escolheram alternativas incorretas, decai a partir da primeira nota não nula como função da nota nesta parte da prova.
Os gráficos relativos às demais questões do componente de Conhecimento Específico constam do Anexo I. 3.3 ANÁLISE DAS QUESTÕES DISCURSIVAS 3.3.1 Componente de Formação Geral A análise dos resultados de desempenho dos estudantes de Tecnologia em Agroindústria, nas duas questões discursivas relativas à Formação Geral, encontra -se na Tabela 3.12 e no Gráfico 3.9.
Observa -se, a partir destes, que as notas médias foram mais baixas nesse conjunto do que no conjunto das objetivas.
A média geral dos Ingressantes nas questões objetivas de Formação Geral foi 52,6 ( Tabela 3.4 ), ao passo que, nas questões discursivas, a média foi 40,6 para esse grupo.
O mesmo ocorreu entre os Concluintes, que tiveram, em Formação Geral, média 54,7 nas questões objetivas e 36,7 nas questões discursivas.
Cabe destacar que, nas questões discursivas, o desempenho dos Concluintes não se evidenciou como superior ao dos Ingressantes.
As notas máximas foram 95,0 para Ingressantes e 77,5 para Concluintes, o que acompanha a constatação de que não houve superioridade dos Concluintes, em relação aos Ingressantes, nesta parte do exame.
Os coeficientes de assimetria indicam assimetria negativa para os Concluintes e para os Ingressantes.
Verifica -se a semelhança das distribuições das notas desses grupos, no Gráfico 3.9, apresentado a seguir.
Na sequência, os resultados verificados para cada uma das questões discursivas de Formação Geral serão apresentados, estabelecendo -se relações com os conteúdos englobados em cada uma delas.
Os comentários da banca corretora a respeito do observado nas respostas dos estudantes, suas impressões e conclusões serão apresentados junto à análise de cada questão.
Cumpre esclarecer que, tendo em vista que as questões discursivas de Formação Geral são padronizadas, ou seja, constam de todas as provas, os comentários da banca são os mesmos para todas as carreiras acadêmicas, sendo direcionados a todos os estudantes que participaram do ENADE/2010.
A seguir, serão analisados os desempenhos de Ingressantes e Concluintes da Área de Tecnologia em Agroindústria nas duas questões discursivas de Formação Geral do ENADE/2010. 3.3.1.1 Análise da Questão Discursiva 9 do Componente de Formação Geral Os dados obtidos a partir das respostas à questão 9 encontram -se na Tabela 3.13 e no Gráfico 3.10.
Nessa questão, os Ingressantes tiveram média 42,0, e os Concluintes, 35,7, resultando com diferença de 6,3 pontos a mais para Ingressantes.
Deixaram a questão em branco 20,9 \% dos Ingressantes e 29,6 \% dos Concluintes.
Mesmo não tendo deixado a questão em branco, obtiveram nota zero na questão 1,5 \% e 1,2 \%, respectivamente.
Entre Ingressantes, a nota mediana foi 45,0 pontos, enquanto Concluintes obtiveram 40,0.
Ambos tiveram maior proporção de notas no intervalo de 50 a 60, com 17,0 \% dos Ingressantes e 16,0 \% dos Concluintes nessa faixa.
A nota máxima, 100,0 pontos, foi alcançada apenas por Ingressantes, ao passo que a de Concluintes foi 85,0. 3.3.1.2 Comentários sobre as respostas à Questão Discursiva 9 No que tange mais precisamente à atualidade do tema proposto e ao desenvolvimento do conteúdo, deve ser ressaltado que se mostrou sintonizado com a realidade política brasileira, bem como com a necessária reflexão sobre o papel na sociedade do profissional formado pelas Instituições de Ensino Superior do país.
Dissertar sobre democracia e ética foi algo apropriado para a época em que a prova foi aplicada, pois coincidiu com o período das eleições, em que tais temas são mais discutidos na vida cotidiana.
Não se pode, entretanto, deixar de se levar em conta a tendência de os estudantes terem limitado a associação do tema “ ética ” ao assunto “ política ”, considerando -se que a ética abarca não apenas a política.
Tal associação atingiu, em alguma medida, o alcance mais preciso das respostas.
As respostas apresentadas refletiram motivação para escrever sobre o tema da questão.
Alguns estudantes chegaram a desenvolver textos \textbf{eloquentes} e engajados.
Em alguns casos, o conceito de democracia, por exemplo, foi bem desenvolvido ao longo do texto, simultaneamente às exemplificações pedidas.
Além disso, muitos estudantes associaram o exercício da ética ao seu futuro ofício, reconhecendo o seu papel como cidadão no desempenho de sua profissão.
Essa associação insinuou, também, o esforço de reconhecimento da noção de ética prática, ainda que de maneira \textbf{codificada}.
Um dos pontos positivos revelou -se no fato de, na maioria das respostas analisadas, haver uma preocupação em definir a sociedade democrática.
Isso facilitou a correção e elevou a média de notas.
Outro aspecto positivo foi a presença pertinente de exposições coerentes de comportamento ético de um futuro profissional.
A articulação dos itens “ a ” e “ b ” da questão produziu, portanto, todo tipo de reflexão acerca de uma postura ética na sociedade, inclusive, corporativa.
A comprovação de que o tema ética tem sido abordado nos cursos superiores foi outro destaque positivo detectado durante as correções.
Assim, além do caráter crítico de que muitas das respostas corretas se revestiram, os relatos constantes sobre desdobramentos da ética comprovaram a abordagem sobre o tema, no mínimo, a ética profissional, nas carreiras universitárias que participaram do ENADE.
A despeito de alguns estudantes se posicionarem contra o exame, seja anulando a questão ou escrevendo algum comentário desabonador sobre o processo, foi observado que significativa parte dos participantes tem consciência político-acadêmica e percebe por que está realizando aquele tipo de prova.
Desse modo, a maioria dos estudantes tinha algo a “ dizer “ sobre o tema da questão, isto é, valia -se da oportunidade para se posicionar a respeito do assunto trazido à \textbf{baila}.
Ainda do ponto de vista do conteúdo desenvolvido nas respostas, observou -se preocupação solidária e consciência cidadã dos estudantes.
Mais um destaque positivo crucial se refere à habilidade textual.
Os estudantes, por exemplo, expressaram atitudes éticas e não éticas e as contrapunham ao conceito de democracia.
Há casos em que os alunos definiam a corrupção política como um exemplo de atitude não ética, enfatizando a real missão desse cargo em uma sociedade democrática : “ atender aos anseios da população e zelar pelo bem comum ”, escreveram vários estudantes.
No tocante à produção de texto dissertativo, à objetividade e ao caráter sintético, muitos estudantes corresponderam ao solicitado, elaborando pequenos textos dissertativos.
Nessas respostas, percebeu -se demonstração de domínio do assunto focado e utilização dos textos de referência como base para a elaboração.
O enunciado da questão 9 determinou que fosse elaborado um texto dissertativo que contivesse a )..., b )... e c )..., fornecendo um encaminhamento para as respostas que conduziu tanto a um resultado positivo quanto a um negativo, a saber : foi produtiva a resposta dos estudantes no sentido de que conseguiram, na maior parte das vezes, atender plenamente aos três quesitos apresentados na questão ; por outro lado, mais especificamente devido ao uso das três letras indicativas ( a, b e c ), houve tendência à elaboração de respostas divididas em itens, em vez de um texto eminentemente dissertativo.
Muitos estudantes apresentaram respostas nesse formato fragmentado para satisfazer a formalidade do roteiro sugerido no comando da questão.
A ideia-chave que deveria permear toda a resposta – democracia – se perdia em fragmentos desconexos entre si.
Se, por um lado, o roteiro do comando da questão facilitou o raciocínio de alguns dos estudantes, por outro, induziu grande parte deles a dois caminhos : i ) o das respostas curtas e objetivas, do tipo “ um profissional de saúde não ético é aquele que conta para as outras pessoas o que foi feito na consulta com o seu paciente ” e ii ) o outro caminho, que é do contraponto entre o conceito de sociedade democrática e a realidade brasileira.
Esse segundo caminho levou vários universitários a produzirem textos bem escritos, distantes, entretanto, do tema exigido.
No que se reporta às especificidades do desenvolvimento do conteúdo, a dificuldade ou falta de competência comunicativa revelada em grande parte das respostas denuncia a pouca exposição dos alunos à leitura de textos escritos de mais excelência.
Tal dificuldade mostrou -se mais evidente na impropriedade no uso de conectores e na incapacidade de elaboração de texto dissertativo, no tratamento de questões mais abstratas.
Muitos estudantes demonstraram não saber em que consistia a ética e confundiram exemplificação com definição : ao tentarem dar exemplos, via de regra, elaboravam uma definição incipiente do que se pedia.
Portanto, a dificuldade dos estudantes em argumentar e articular as ideias com coesão e coerência ainda são os grandes “ vilões ” desse processo avaliativo.
O excesso de fragmentos e textos prévios acerca do tema permitiu que a elaboração das respostas se resumisse a uma colagem ou paráfrase mal feita dos materiais disponíveis.
Em vez de subsidiar as respostas, esse material acabou por influenciar negativamente o desempenho dos estudantes.
Desse modo, verificou -se um nível de letramento ainda significativamente baixo, visto que a maioria apresentou uma visão bastante reducionista e simplificada sobre o tema em pauta.
De um modo geral, ao conceituarem ética e sociedade democrática, os estudantes demonstraram desconhecer o limite entre os dois conceitos, sem aludirem à interseção de ambos.
Foi expressiva a quantidade de redações nas quais questionavam a democracia pelo simples fato de estarem fazendo, a fortiori, aquela prova especificamente.
Em muitas das respostas, os estudantes se diziam incomodados pelo fato de estarem obrigados a realizar o exame, demonstrando sua insatisfação.
Por diversas vezes, os estudantes não foram capazes de sintetizar e de selecionar seu discurso o suficiente para atender ao que era pedido, o que ocasionou a falta de, pelo menos, um dos aspectos da resposta.
Foi constatada grande parte de respostas com problemas graves em termos textuais e gramaticais, em relação ao que se esperaria de pessoas que estão vinculadas ao Ensino ( Técnico ) Superior.
Alguns estudantes sequer conseguiam elaborar uma frase coerente e coesa, apresentando respostas que poderiam ser classificadas como de analfabetismo funcional.
De maneira ampla, manifestou -se da parte dos estudantes uma grande dificuldade no trato da língua escrita, dificuldade essa que se traduz, sobretudo, por um desrespeito bastante frequente às regras relativas : i ) à ortografia ( por exemplo : “ proficionais ”, “ infligir as leis ”, “ corrupitos ”, “ democrácia ”, “ opitar ”, “ manter o cizilo proficionau ” ), troca de “ sociedade democrática ” por “ sociedade demográfica ” e de “ ética ” por “ étnica ” ; ii ) à pontuação ; iii ) à concordância verbal ( “ os político deve (... ) eles pode ” ) ; iv ) ao modo de organização textual ( sem distinção entre o que seja dissertativo, argumentativo ou narrativo ).
Frisa -se que esse comentário não se restringe ao que se costuma chamar de “ língua padrão ”, mas ao que poderia ser chamado de “ excelência textual ”, que não se mede exclusiva ou enfaticamente pela correção gramatical, mas pela competência, objetividade, clareza e, em alguma medida, pela criatividade no ato de redigir.
Por fim, destaca -se um dos problemas dignos de nota na formulação da questão.
A transcrição do Dicionário Houaiss apresentava defeitos.
Por exemplo, a definição 5 de “ democracia ” indicava que esse é o tipo de “ governo que acata a vontade da maioria, embora respeitando os direitos e a livre expressão das minorias ”.
O emprego de um valor concessivo na oração adverbial, por oposição a um valor aditivo, contradiz as concepções efetivamente democráticas apresentadas de 1 a 4.
Repercussões desse valor concessivo, mal empregado, foram encontradas em incontáveis respostas, as quais diziam ser típico de um governo democrático, seguir a vontade da maioria, “ apesar ” de respeitar as minorias.
Outro desses problemas dignos de nota refere -se à indicação do dicionário quanto a “ governo ”, ao passo que o comando focava a ideia mais ampla de “ sociedade democrática ”. 3.3.1.3 Análise da Questão Discursiva 10 do Componente de Formação Geral O desempenho dos estudantes na questão 10 foi inferior ao obtido na questão de número 9, quanto à nota média de Ingressantes, que foi 39,1.
A média de Concluintes foi um pouco mais elevada do que na questão anterior, atingindo 37,7.
A nota máxima foi alcançada em ambos os grupos, sendo que os percentuais de alunos que deixaram a questão em branco foram de 22,3 \% para Ingressantes e 25,9 \% para Concluintes.
Houve maior percentual de alunos, considerando -se a distribuição de notas não nulas, no intervalo entre 40 e 50 pontos, tanto para Ingressantes quanto para Concluintes, em torno de 18,0 \%.
Não se manifestou superioridade de Concluintes nas diferentes faixas de notas posteriores a esse intervalo, com equilíbrio no desempenho dos dois grupos, só ligeiramente superior no intervalo de notas entre 60 e 70.
Essas informações estão ilustradas na Tabela 3.14 e no Gráfico 3.11. 3.3.1.4 Comentários sobre as respostas à Questão Discursiva 10 Inicialmente, será apresentada uma panorâmica sobre as respostas, com destaque para os aspectos positivos observados.
Em relação ao desenvolvimento do tema e à vinculação do mesmo à vida cotidiana dos estudantes que participaram do exame, a visão positiva em relação à Educação de Jovens e Adultos ( EJA ) reforça a ideia de que a educação pode trazer crescimento não somente profissional, mas, sobretudo, para ampliar os \textbf{horizontes} da cidadania.
Foram frequentes as respostas que apontavam a educação como a melhor saída, senão a única, para o resgate da cidadania.
A atuação dos governos, em especial o governo federal, foi um ponto destacado nas respostas, as quais sinalizavam para a necessidade de as instâncias governamentais assumirem um papel cada vez mais ativo na consolidação da educação – pública, sobretudo – no país.
As políticas públicas focadas na educação foram, então, um item bastante frequente nas respostas, seja para enaltecer, criticar ou requisitar a ação governamental.
Outro ponto positivo destacado diz respeito à consciência ou ao reconhecimento da precariedade de condições de trabalho dos professores, em geral, no Brasil e, em particular, daqueles ocupados com turmas de jovens e adultos.
Nesse sentido, foram inúmeras as propostas apresentadas na segunda parte das respostas com relação à imperiosidade de melhoria salarial para os docentes, e à melhoria das condições gerais de trabalho desses profissionais do ensino.
A associação direta entre valorização do professor e sucesso escolar dos estudantes foi um ponto complementar a essa preocupação com os docentes.
Os estudantes demonstraram uma visão de vínculo no processo de ensino-aprendizagem, no qual alunos e professores são partes inalienáveis entre si.
A questão da repetência e da evasão escolar, referenciadas no suporte da questão, foram então, abordadas pelos alunos por meio do viés dessa valorização docente.
Ainda que os temas específicos da repetência e da evasão não fossem contemplados a contento na maior parte das respostas, ao menos, nessa associação com a qualificação e a valorização dos professores, algo foi trabalhado analiticamente pelos estudantes.
Num plano mais abrangente, houve reconhecimento de que, no processo ensino-aprendizagem, não há atores isolados : alunos, profissionais de educação, governos e a própria sociedade como \textbf{corresponsáveis} pelo sucesso ou fracasso de tal processo.
A atualidade do tema e a sua explícita vinculação com o cotidiano de muitos estudantes que realizavam o exame fizeram com que os respondentes demonstrassem um visível grau de envolvimento com o tema proposto.
Ou seja, foi possível detectar um empenho notório dos estudantes ao redigir suas respostas com seriedade e comprometimento relativo à delicada problemática dos milhões de estudantes trabalhadores do país.
Utilizar todo o espaço destinado à resposta, procurando aproveitar ao máximo a oportunidade de emitir um argumento ou uma opinião revelou -se como uma das evidências do aludido empenho.
A EJA também foi reconhecida como um dos caminhos concretos para o acesso à Universidade, no médio prazo.
As respostas se voltaram, muitas vezes, à condição de depoimentos pessoais sobre a conciliação de trabalho e estudo.
Nessa direção, ao assumirem o papel de depoentes, os estudantes expuseram suas histórias de vida, suas expectativas e percalços, suas \textbf{esperanças} e \textbf{entraves} cotidianos, seus anseios mais ambiciosos e problemas estruturais, que poderiam se resolver por meio da educação.
Nesse inciso, ficava clara a relação que se formava entre educação e mobilidade social, mediatizada por iniciativas governamentais como a EJA.
Uma sugestão plausível e válida para a solução dos problemas com a EJA no país reportou -se à necessidade de maior participação do público alvo na definição de estratégias para essa modalidade de ensino.
Nesse detalhe, era notória a associação das preocupações relativas à questão 9, que trata do problema da democracia e da cidadania, com o problema específico da educação escolar.
Houve um número significativo de sugestões ou propostas para se combater o problema da repetência e da evasão, destacando -se melhores salários para os professores e aulas mais dinâmicas.
A lista dessas sugestões, contudo, era bem mais ampla e diversificada, incluindo ainda : a redução da jornada de trabalho dos jovens e adultos que participam da EJA ; melhoria de infraestrutura nas instituições responsáveis pela EJA ; incentivo financeiro aos estudantes trabalhadores por meio de um sistema de bolsas de estudo ; realização de campanhas motivacionais para o retorno aos estudos ; estímulo à participação das empresas, por meio da solicitação de recursos materiais, espaço físico para as salas, auxílio financeiro ; incremento dos cursos profissionalizantes no país ; e ampliação da oferta de ensino a distância em todo o território nacional.
Muitos estudantes abordaram, com propriedade, a associação entre um mercado de trabalho cada vez mais competitivo e a necessidade de qualificação profissional, inclusive com uma clara percepção de que os salários são tanto maiores quanto maior a quantidade de anos de estudo e/ou a qualidade da formação.
Nesse tópico, retorna -se ao papel crucial da EJA como viabilizadora dessa formação.
Em continuidade à análise, serão apresentados aspectos negativos que foram notados com frequência nas respostas à questão 10.
De um ponto de vista mais global, a dificuldade repousa no entendimento da questão, isto é, na decifração do comando.
Muitos fugiram completamente do assunto proposto, ainda que o enunciado da questão fosse farto em dados específicos ao tema.
Em vez de responderem diretamente quais eram os problemas enfrentados por aqueles que têm de trabalhar e estudar ao mesmo tempo, muitos estudantes optavam por falar das dificuldades que existem na educação brasileira, de forma estendida e pasteurizada.
Vários estudantes não conseguiam, em suas respostas, relacionar as dificuldades de conciliação da escola com o trabalho, a não ser, como já mencionado, nos casos de depoimentos pessoais, relatos de experiência, ou seja, traços de uma biografia.
Assim, eram respostas cujo teor era apenas o de registrar o abandono dos estudos por parte de vários jovens que têm de trabalhar.
O tom das respostas, nesses casos, era meramente de uma descrição e não de uma análise, muito menos de um elemento que compusesse uma argumentação.
A dificuldade de compreensão da questão se expressava nas respostas que, em lugar de mostrarem os problemas de conciliação entre escola e trabalho, prendiam -se à ideia de “ associação ” – termo que consta do enunciado da questão – entre escola e trabalho.
Muitos alunos relatavam as consequências positivas de uma boa educação no momento em que se busca bons postos de trabalho.
Tratava -se de uma fuga parcial ao comando da questão, evidenciando má compreensão do mesmo.
Uma dificuldade facilmente percebida nas respostas refere -se à tendência de os estudantes tratarem dos problemas relativos à escola e, muito separadamente, daqueles problemas relativos ao trabalho, reconhecendo uma espécie de sinergia que impediria de visualizá -los em sua individualidade, ao mesmo tempo, que em sua complementaridade.
Isso significa que os estudantes distinguiam por mero isolamento, ou seja, eles não conseguiam distinguir e articular e, sobretudo, reconhecer as convocações mútuas entre os problemas da escola e do trabalho, no caso explícito dos estudantes trabalhadores.
Essa dificuldade de identificar, associar e analisar aspectos distintos, até mesmo concorrentes, antagônicos e complementares, denuncia a fragilidade de um pensamento simplificador que a qualidade das respostas revelou ser ainda hegemônico nas universidades brasileiras.
Notou -se a presença de preconceitos na análise dos fatores que levam ao abandono dos estudos.
Muitas respostas apontavam como causa desse abandono e da repetência a preguiça, a falta de perseverança e o cansaço dos estudantes.
Nesse tópico específico, fica evidente a responsabilização do indivíduo pelos percalços e fracassos de sua vida, descontextualizando -o de uma sociedade que, em seu movimento de conjunto, lhe impõe condicionantes.
A precariedade do uso da língua pátria de modo culto foi uma constante, conforme mencionado na análise da questão 9.
As impropriedades lexicais e a falta de elementos coesores que articulassem as partes constitutivas das respostas foram um incômodo constante.
Nesse sentido, agrava -se a situação da baixa qualidade textual em decorrência complementar da construção de textos meramente descritivos e não argumentativos, como é exigido no suporte da questão.
Isso demonstra a dificuldade dos estudantes universitários na compreensão de que uma argumentação requer, imperiosamente, justificações para os apontamentos apresentados.
Os textos das respostas eram, em sua maioria, opinativos, sem qualquer compromisso com elementos explicativos ou com a busca de justificativas. 3.3.2 Componente de Conhecimento Específico Na parte da prova relativa ao componente de Conhecimento Específico, observa -se que a diferença de desempenho foi mais acentuada entre Ingressantes e Concluintes.
A Tabela 3.15 contém dados elucidativos, na qual é possível observar que, no grupo dos Ingressantes, a nota média 11,1 é inferior à alcançada por Concluintes, 19,0.
Essas médias foram menores do que as obtidas no conjunto das questões objetivas, nas quais Ingressantes alcançaram nota média 45,3, e Concluintes, 52,3 ( Tabela 3.8 ).
O Gráfico 3.12 apresenta a distribuição de notas nos diferentes intervalos.
Observa -se que tanto Ingressantes quanto Concluintes atingiram proporções mais elevadas na faixa entre zero e dez pontos, com 57,3 \% e 38,3 \%, respectivamente.
Acima dessa faixa, houve concentração significativa de Concluintes no intervalo entre 20 e 30, com 23,5 \% dos estudantes.
A distribuição das notas dos Ingressantes é decrescente com as faixas.
A análise de cada uma destas questões será feita a seguir. ∗ 3.3.2.2 Comentários sobre as respostas à Questão Discursiva 38 Em parte significativa das respostas elaboradas pelos estudantes, os termos Análise de Pontos Críticos de Controle ( APPCC ) e Pontos de Controle ( PC ) foram ∗ A análise completa, com apresentação de trechos de respostas de estudantes, pode ser encontrada no documento “ Relatório da Correção das Questões Discursivas ”, elaborado pelos docentes que participaram do processo de correção das mesmas. 3.3.2.1 Análise da Questão Discursiva 38 do Componente de Conhecimento Específico Na questão 38, cujos resultados aferidos encontram -se descritos na Tabela 3.16, o grupo dos Ingressantes obteve nota média 11,7, enquanto a média dos Concluintes foi 22,7.
Deixaram a questão em branco 37,9 \% dos Ingressantes e 30,9 \% dos Concluintes.
As notas dos Concluintes foram superiores às notas dos Ingressantes a partir do intervalo 30 a 40 pontos.
A nota máxima, para Ingressantes, foi 75,0, enquanto Concluintes obtiveram 85,0. 3.3.2.2 Comentários sobre as respostas à Questão Discursiva 38 Em parte significativa das respostas elaboradas pelos estudantes, os termos Análise de Pontos Críticos de Controle ( APPCC ) e Pontos de Controle ( PC ) foram utilizados como se fossem a mesma coisa.
Uma minoria se referiu à presença de metais pesados na água de irrigação, e praticamente nenhum se referiu aos perigos físicos, no item “ a ”.
Verificou -se que, na maioria das respostas, foi lembrada a importância do uso correto de substância sanificantes para a lavagem dos produtos.
No item “ b ”, os estudantes simplesmente citavam alguma medida preventiva, sem relacioná -la aos perigos. 3.3.2.3 Análise da Questão Discursiva 39 do Componente de Conhecimento Específico A Tabela 3.17 contém as informações relativas à questão 39 do conjunto de questões do componente de Conhecimento Específico.
O desempenho dos dois grupos – Ingressantes e Concluintes – nessa questão, foi um pouco superior ao desempenho na questão 38.
As médias foram 14,7 no grupo dos Ingressantes e 23,5 no grupo dos Concluintes.
A proporção de estudantes que deixaram a questão em branco foi semelhantes a observada na questão 38, sendo esse percentual de 36,4 \% entre os Ingressantes e 32,1 \% entre os Concluintes.
A nota máxima atingida por Ingressantes ( 80,0 ) foi superior a dos Concluintes ( 70,0 ).
Ingressantes tiveram maior proporção de notas zero, mesmo considerando -se aqueles que não deixaram a questão em branco : 14,6 \%.
O intervalo de 30 a 40 pontos corresponde às notas não nulas de maior frequência entre Concluintes com 13,6 \%.
Enter Ingressantes a faixa modal foi de 20 a 30 pontos com 12,1 \% dos estudantes do grupo.
As distribuições das notas entre os estudantes constam do Gráfico 3.14. 3.3.2.4 Comentários sobre as respostas à Questão Discursiva 39 No item “ a ”, boa parte dos estudantes se referiu corretamente ao corte da juçara como fator de extinção de espécie, porém, em uma quantidade significativa de respostas, afirmaram que a juçara apresentava maior risco de extinção por estar presente em várias áreas do país, enquanto o açaí se restringia à região amazônica.
No item “ b ”, dentre as alternativas tecnológicas mais citadas, estava o replantio.
Já no item “ c ”, os estudantes lembraram a importância do controle do C. botulinum, do pH e da temperatura como medidas de controle de qualidade para o produto.
Quanto ao item “ d ”, a grande maioria citou a utilização dos resíduos como matéria-prima para artesanato. 3.3.2.5 Análise da Questão Discursiva 40 do Componente de Conhecimento Específico A Tabela 3.18 contém as informações relativas à questão 40 do conjunto de questões do componente de Conhecimento Específico.
O desempenho dos dois agrupamentos de estudantes, nessa questão, foi inferior ao ser comparado às questões 38 e 39.
A nota média de Ingressantes foi 6,7 ; já a de Concluintes, 10,8.
A nota mediana total foi zero.
As notas máximas obtidas foram 70,0 para Ingressantes e 65,0 para Concluintes.
A questão foi deixada em branco por 55,8 \% dos Ingressantes e por 44,4 \% dos Concluintes.
A proporção de estudantes que não deixaram a questão em branco mas receberam nota zero foi 12,6 \% entre Ingressantes e 13,6 \% entre Concluintes.
O intervalo de notas não nulas no qual foi observada maior frequência corresponde à faixa que vai de zero a 10 pontos, com 12,6 \% dos Ingressantes e 12,3 \% dos Concluintes. o que demonstra ter sido esta a questão de desempenho mais baixo.
A distribuição de notas nos dois agrupamentos pode ser observada no Gráfico 3.15. 3.3.2.6 Comentários sobre as respostas à Questão Discursiva 40 No item “ a ”, os estudantes, de maneira geral, se manifestaram de forma correta sobre a produção de cada fruta em diferentes épocas do ano para dimensionar os equipamentos e os gastos com insumos, bem como em relação à capacidade de processar da indústria.
Pouquíssimos, entretanto, apresentaram uma real noção de quais dados são necessários para o dimensionamento da planta de produção solicitada na questão.
No item “ b ”, muitos estudantes citaram corretamente os equipamentos, porém poucos citaram os quatro de forma correta ( o comum foi a citação de 2 ou 3 ).
Nenhum graduando conseguiu atingir a nota máxima final ( 100,0 ), principalmente devido ao baixo índice de acertos do item “ a ”.
O item “ b ”, por outro lado, foi o que apresentou maior índice de acertos.
CONSIDERAÇÕES FINAIS De um modo geral, a despeito das questões terem sido as últimas da prova, os estudantes tiveram um rendimento médio aceitável.
Alguns responderam de forma correta e objetiva.
Foi notória a ausência de conhecimento básico de disciplinas correlatas, por exemplo, na questão que relacionava Açaí e Juçara, de certo modo tema atual no Brasil.
Erros de localização geográfica e de conhecimento de nossos biomas foram comuns.
Ficou faltando demonstração de conhecimento básico em relação aos cuidados necessários para o dimensionamento de uma indústria.
Espera -se que as informações aqui contidas possam contribuir para a melhor qualificação dos estudantes que estão cursando Tecnologia em Agroindústria, especialmente para os que hão de vir, que, certamente, usufruirão do aprimoramento do trabalho realizado pelas instituições de ensino, advindo de avaliações como o ENADE.

% matched lemmas: baila, codificar, corresponsável, eloquente, entrave, esperança, horizonte
\end{textsample}
